% Chinese reading copy of the documented School of Computing long-thesis adaptation.
\usepackage{fontspec}
\usepackage{xeCJK}
\setmainfont{lmroman10-regular.otf}[
  BoldFont=lmroman10-bold.otf,
  ItalicFont=lmroman10-italic.otf,
  BoldItalicFont=lmroman10-bolditalic.otf]
\setsansfont{lmsans10-regular.otf}[
  BoldFont=lmsans10-bold.otf,
  ItalicFont=lmsans10-oblique.otf,
  BoldItalicFont=lmsans10-boldoblique.otf]
\setmonofont{lmmono10-regular.otf}[
  BoldFont=lmmonolt10-bold.otf,
  ItalicFont=lmmono10-italic.otf,
  BoldItalicFont=lmmonolt10-boldoblique.otf]
\setCJKmainfont{FandolSong-Regular.otf}[
  BoldFont=FandolSong-Bold.otf,
  ItalicFont=FandolKai-Regular.otf]
\setCJKsansfont{FandolHei-Regular.otf}[BoldFont=FandolHei-Bold.otf]
\setCJKmonofont{FandolFang-Regular.otf}
\XeTeXlinebreaklocale "zh"
\XeTeXlinebreakskip=0pt plus 1pt
\usepackage{microtype}
\usepackage{geometry}
\usepackage{setspace}
\usepackage{xcolor}
\definecolor{ThesisNavy}{HTML}{18344A}
\definecolor{DraftColour}{HTML}{8A4515}
\usepackage{graphicx}
\graphicspath{{figures/}}
\usepackage{amsmath,amssymb,amsthm}
\usepackage{booktabs,array,tabularx,longtable}
\usepackage{caption,subcaption}
\captionsetup{font=small,labelfont=bf}
\captionsetup[table]{position=top}
\usepackage{enumitem}
\setlist{itemsep=0.2em,topsep=0.35em}
\usepackage{listings}
\lstset{basicstyle=\small\ttfamily,breaklines=true,
  columns=fullflexible,keepspaces=true,showstringspaces=false,
  frame=single,rulecolor=\color{black!20},tabsize=2}
\usepackage{tikz}
\usetikzlibrary{arrows.meta,positioning,decorations.pathreplacing}
\usepackage[numbers,sort&compress]{natbib}
\usepackage[automark,headsepline]{scrlayer-scrpage}
% KOMA-Script long-thesis adaptation; see FORMAT_NOTES.txt for the ANU source.
% Independently configured with public KOMA interfaces; no third-party template source is redistributed.
% Keep the existing font size and margins. Normal leading follows the school template.
\geometry{a4paper,left=27mm,right=27mm,top=28mm,bottom=28mm,
  headheight=18pt,headsep=7mm,footskip=12mm}
\setstretch{1.0}
% Classic thesis paragraphs: indent instead of a half-line blank between them.
\setparsizes{1.2em}{0pt plus .5pt}{0pt plus 1fil}
% Preserve the existing, scoped roman math subscripts without rewriting equations.
\DeclareOldFontCommand{\rm}{\normalfont\rmfamily}{\mathrm}

% Mathematical statements and proofs share the body-text margins. Their
% headings occupy a separate line; upright statements distinguish prose from
% the mathematical symbols without switching whole paragraphs to italics.
\newtheoremstyle{thesis}
  {6pt plus 1pt minus 1pt}{4pt plus 1pt minus 1pt}
  {\normalfont\setlength{\parindent}{0pt}}{0pt}
  {\normalfont\bfseries}{.}{\newline}{}
\theoremstyle{thesis}
\makeatletter
\renewenvironment{proof}[1][\proofname]{%
  \par\pushQED{\qed}\normalfont
  \setlength{\parindent}{0pt}\setlength{\listparindent}{0pt}%
  \setlength{\parskip}{2pt plus .5pt}%
  \topsep4\p@\@plus\p@\@minus\p@\relax
  \trivlist
  \item[\hskip\labelsep\normalfont\bfseries #1\@addpunct{.}]%
  \leavevmode\par\nobreak\noindent\ignorespaces
}{%
  \popQED\endtrivlist\@endpefalse
}
\makeatother
% Each current statement/proof pair fits comfortably on one page. Keep the
% short pairs together so floats and page breaks cannot separate their parts.
\newenvironment{resultblock}{%
  \par\addvspace{6pt}\noindent\begin{minipage}[t]{\linewidth}%
}{%
  \end{minipage}\par\addvspace{6pt}%
}

\KOMAoptions{chapterprefix=false,headings=normal,numbers=noendperiod}
\setkomafont{disposition}{\normalfont\bfseries}
\setkomafont{chapter}{\normalfont\Large\bfseries}
\setkomafont{chapterprefix}{\normalfont\large\bfseries}
\RedeclareSectionCommand[beforeskip=1.5ex plus .3ex minus .2ex,
  afterskip=.8ex plus .2ex]{chapter}
\RedeclareSectionCommand[beforeskip=1.8ex plus .4ex minus .2ex,
  afterskip=.7ex plus .2ex]{section}
\RedeclareSectionCommand[beforeskip=1.4ex plus .3ex minus .2ex,
  afterskip=.6ex plus .2ex]{subsection}
\RedeclareSectionCommand[runin=true,beforeskip=.9ex plus .3ex minus .1ex,
  afterskip=.75em]{paragraph}

\clearpairofpagestyles
\automark[chapter]{chapter}
\ihead{\headmark}
\cfoot[\pagemark]{\pagemark}
\setkomafont{pageheadfoot}{\normalfont\footnotesize\color{black!65}}
\setkomafont{pagenumber}{\normalfont\small\color{black}}
\pagestyle{scrheadings}

\captionsetup{font={small,stretch=1.0},labelfont=bf,labelsep=period,
  justification=raggedright,singlelinecheck=false,skip=5pt}
\captionsetup[table]{position=top,skip=4pt}
\setlength{\textfloatsep}{12pt plus 2pt minus 2pt}
\setlength{\floatsep}{10pt plus 2pt minus 2pt}
\setlength{\intextsep}{10pt plus 2pt minus 2pt}
% Keep neighbouring floats together on float-only pages.
\makeatletter
\setlength{\@fptop}{0pt}
\setlength{\@fpsep}{12pt plus 2pt minus 2pt}
\setlength{\@fpbot}{0pt plus 1fil}
\makeatother
\setlength{\abovecaptionskip}{5pt}
\setlength{\belowcaptionskip}{0pt}
\setlist{itemsep=.15em,topsep=.35em,parsep=0pt}
\setlength{\bibsep}{.3em}
\setlength{\abovedisplayskip}{8pt plus 2pt minus 2pt}
\setlength{\belowdisplayskip}{8pt plus 2pt minus 2pt}
\setlength{\abovedisplayshortskip}{4pt plus 1pt}
\setlength{\belowdisplayshortskip}{5pt plus 1pt minus 1pt}
\AtBeginDocument{%
  \setlength{\abovedisplayskip}{8pt plus 2pt minus 2pt}%
  \setlength{\belowdisplayskip}{8pt plus 2pt minus 2pt}%
  \setlength{\abovedisplayshortskip}{4pt plus 1pt}%
  \setlength{\belowdisplayshortskip}{5pt plus 1pt minus 1pt}}

% Short preliminary sections may share a page; numbered chapters stay separate.
\newcommand{\frontsection}[1]{%
  \section*{#1}\phantomsection\addcontentsline{toc}{chapter}{#1}%
  \markboth{#1}{#1}}

% Native KOMA contents/list styles and appendix terminology.
\DeclareTOCStyleEntry[beforeskip=.5em plus .1em]{chapter}{chapter}
\DeclareTOCStyleEntries[pagenumberwidth=2.5em]{chapter}{chapter}
\DeclareTOCStyleEntry[indent=0pt,numwidth=2.8em]{tocline}{figure}
\DeclareTOCStyleEntry[indent=0pt,numwidth=2.8em]{tocline}{table}
\AddToHook{cmd/appendix/after}{\crefalias{chapter}{appendix}}

% Chinese characters need slightly more leading than the Latin reading copy.
\setstretch{1.25}
\setparsizes{2em}{0pt plus .5pt}{0pt plus 1fil}
\renewcommand*{\chapterformat}{第\thechapter 章}
\renewcommand*{\chaptermarkformat}{第\thechapter 章\enskip}
\AddToHook{cmd/appendix/after}{%
  \renewcommand*{\chapterformat}{附录~\thechapter}%
  \renewcommand*{\chaptermarkformat}{附录~\thechapter\enskip}}
\setcounter{secnumdepth}{2}
\setcounter{tocdepth}{1}
\setlength{\emergencystretch}{2em}
\widowpenalty=10000
\clubpenalty=10000
\raggedbottom
% Match the English thesis's result-figure placement.
\renewcommand{\topfraction}{0.90}
\renewcommand{\bottomfraction}{0.85}
\renewcommand{\textfraction}{0.10}
\renewcommand{\floatpagefraction}{0.75}
\setcounter{topnumber}{3}
\setcounter{totalnumber}{4}
\makeatletter
\setlength{\@fptop}{0pt}
\makeatother
\usepackage[hyphens]{url}
\usepackage[unicode,hidelinks,pdfusetitle]{hyperref}
\hypersetup{pdftitle={\thesistitle},pdfauthor={\thesisauthor},
  pdfsubject={COMP8800 research thesis},bookmarksnumbered=true}
\usepackage[nameinlink,noabbrev]{cleveref}

\newtheorem{proposition}{命题}[chapter]
\newtheorem{lemma}{引理}[chapter]
\newtheorem{corollary}[proposition]{推论}
\crefname{corollary}{推论}{推论}
\Crefname{corollary}{推论}{推论}
\theoremstyle{definition}
\newtheorem{definition}[proposition]{定义}
\renewcommand{\proofname}{证明}

\newcommand{\frontchapter}[1]{%
  \chapter*{#1}\phantomsection\addcontentsline{toc}{chapter}{#1}
  \markboth{#1}{#1}}
\newcommand{\draftnote}[1]{%
  \ifdraftthesis
    \par\begingroup\small\color{DraftColour}
    \noindent\textbf{待撰写：} #1\par\endgroup\medskip
  \else
    \PackageError{comp8800-template}{Unresolved writing placeholder}
      {Replace all draftnote commands before a final build.}
  \fi}
\newcommand{\figureplaceholder}[3]{%
  \ifdraftthesis
    \begin{figure}[tbp]\centering
      \fbox{\parbox[c][35mm][c]{0.87\linewidth}{\centering\small #1}}
      \caption{\textbf{Draft placeholder.} #2}\label{#3}
    \end{figure}
  \else
    \PackageError{comp8800-template}{Unresolved figure placeholder}
      {Replace figureplaceholder with the final figure.}
  \fi}
\DeclareMathOperator{\Recall}{Recall}
\newcommand{\rawdist}{D_{\mathrm{raw}}}
\newcommand{\primarydist}{D_{\mathrm{P}}}
\newcommand{\refineddist}{D_{\mathrm{H}}}
\newcommand{\glut}{g_{\mathrm{LUT}}}
\newcommand{\gscan}{g_{\mathrm{scan}}}
\newcommand{\lrefine}{\ell_{\mathrm{ref}}}

% Chinese reading-copy names; labels and numbering follow the English draft.
\renewcommand{\bibfont}{\raggedright}
\crefname{figure}{图}{图}
\Crefname{figure}{图}{图}
\crefname{table}{表}{表}
\Crefname{table}{表}{表}
\crefname{proposition}{命题}{命题}
\Crefname{proposition}{命题}{命题}
\crefname{lemma}{引理}{引理}
\Crefname{lemma}{引理}{引理}
\crefname{appendix}{附录}{附录}
\Crefname{appendix}{附录}{附录}
\crefformat{appendix}{附录~#2#1#3}
\Crefformat{appendix}{附录~#2#1#3}
\crefformat{chapter}{第#2#1#3章}
\Crefformat{chapter}{第#2#1#3章}
\crefformat{section}{第#2#1#3节}
\Crefformat{section}{第#2#1#3节}
\crefformat{equation}{式~(#2#1#3)}
\Crefformat{equation}{式~(#2#1#3)}
\AtBeginDocument{%
  \renewcommand{\contentsname}{目录}%
  \renewcommand{\listfigurename}{插图目录}%
  \renewcommand{\listtablename}{表格目录}%
  \renewcommand{\figurename}{图}%
  \renewcommand{\tablename}{表}%
  \renewcommand{\bibname}{参考文献}%
  \renewcommand{\appendixname}{附录}%
  \renewcommand{\crefpairconjunction}{与}%
  \renewcommand{\crefmiddleconjunction}{、}%
  \renewcommand{\creflastconjunction}{与}%
  \renewcommand{\crefpairgroupconjunction}{与}%
  \renewcommand{\crefmiddlegroupconjunction}{、}%
  \renewcommand{\creflastgroupconjunction}{与}%
}
